% Zone „Das kennst du schon“ – aus bank/zufallsgroessen-und-verteilungen/zone.jsonl (zusammenbau v0.1)
\einheitenkopf[zone][Kennst du schon]{Das kennst du schon}
\setcounter{aufgabe}{0}

%% TODO Titel: Ich-kann-Satz fehlt in der Bank (Platzhalter: Kettenname) – Nr. 1
%% TODO Anweisung: ein Satz über der ersten Teilaufgabe fehlt in der Bank – Nr. 1
\begin{aufgabe}{Ergebnismengen zweistufiger Experimente aufzählen und günstige Fälle abzählen (Paare, Reihenfolge beachten)}
\begin{teile}
\teil Eine Münze wird zweimal geworfen (K Kopf, Z Zahl) – wie viele Ergebnisse?
\teil Zwei Würfel werden geworfen – wie viele der $36$ Paare haben die Summe $5$?
\end{teile}
\end{aufgabe}

%% TODO Titel: Ich-kann-Satz fehlt in der Bank (Platzhalter: Kettenname) – Nr. 2
%% TODO Anweisung: ein Satz über der ersten Teilaufgabe fehlt in der Bank – Nr. 2
\begin{aufgabe}{Brüche addieren und zur Summe eins ergänzen}
\begin{teile}
\teil $\frac{1}{5} + \frac{2}{5} + \frac{1}{5}$ – was fehlt zu $1$?
\teil $\frac{1}{4} + \frac{1}{3}$ – was fehlt zu $1$?
\end{teile}
\end{aufgabe}

%% TODO Titel: Ich-kann-Satz fehlt in der Bank (Platzhalter: Kettenname) – Nr. 3
%% TODO Anweisung: ein Satz über der ersten Teilaufgabe fehlt in der Bank – Nr. 3
\begin{aufgabe}{Säulendiagramme lesen (Höhen, Symmetrie, Verhältnisse)}
\begin{teile}
\teil Lies im Diagramm ab – wie hoch ist die Säule Di?

\saeulenab[ymax=10,ystep=2,yfein=1,ylabel=Anzahl]{Mo/4,Di/8,Mi/6}
\teil Lies im Diagramm ab – wievielmal so hoch ist die Säule B wie die Säule A?

\saeulenab[ymax=10,ystep=2,yfein=1,ylabel=Anzahl]{A/3,B/9,C/6}
\end{teile}
\end{aufgabe}

%% TODO Titel: Ich-kann-Satz fehlt in der Bank (Platzhalter: Kettenname) – Nr. 4
%% TODO Anweisung: ein Satz über der ersten Teilaufgabe fehlt in der Bank – Nr. 4
\begin{aufgabe}{Ergebnismengen zweistufiger Experimente aufzählen und günstige Fälle abzählen (Paare, Reihenfolge beachten) – Fehler finden}
\begin{teile}
\teil Zwei Würfel werden geworfen, Summe $4$ – wo ist der Fehler? \rechnung{\text{günstig: } &1 \text{ und } 3,\ 2 \text{ und } 2 \\ P &= \tfrac{2}{36}}
\end{teile}
\end{aufgabe}

%% TODO Titel: Ich-kann-Satz fehlt in der Bank (Platzhalter: Kettenname) – Nr. 5
%% TODO Anweisung: ein Satz über der ersten Teilaufgabe fehlt in der Bank – Nr. 5
\begin{aufgabe}{Ergebnismengen zweistufiger Experimente aufzählen und günstige Fälle abzählen (Paare, Reihenfolge beachten) – selbst rechnen}
\begin{teile}
\teil Zwei Würfel werden geworfen – Wahrscheinlichkeit der Summe $9$?
\end{teile}
\end{aufgabe}
